% Sections 5-8 for the CVPR 2027 paper. Numbers come only from PDF version (11).
% Paste over Sec. 5 onward. Keep the existing tables and figures.
% \CHECK = conflicts with another part of the PDF, or needs confirming.
% \TODO  = needs data that is not in the PDF. No number has been filled in.
% Label names are guesses: tab:replay, tab:selection, tab:storage, fig:q2a, fig:q3.
\providecommand{\CHECK}[1]{\textcolor{orange}{#1}}

%%%%%%%%%%%%%%%%%%%%%%%%%%%%%%%%%%%%%%%%%%%%%%%%%%%%%%%%%%%%%%%%%%%%%%%%%
\section{Results}
\label{sec:results}

\subsection{Main results: replay + distillation (\Q{1})}
\label{sec:res-q1}

\paragraph{Baselines.}
We compare against ZSCL~\cite{zheng2023zscl}, which distills on reference data
without replay and is our direct predecessor; MoE-Adapters~\cite{yu2024moe},
which routes inputs through task-specific experts and keeps no exemplars; and
GIFT~\cite{gift2025} and LoRA-Loop~\cite{loraloop2025}, which replay images
synthesized by a ${\approx}900$M-parameter diffusion model. ZAF~\cite{gao2024zaf},
C-CLIP~\cite{liu2025cclip}, and SD-LoRA~\cite{wu2025sdlora} use different
protocols and are not directly comparable on this benchmark.

\paragraph{Comparison with prior work.}
On Order~I, image replay with distillation reaches $86.16$ Last, 0.16 above
GIFT and LoRA-Loop (86.0) and 2.56 above ZSCL (83.6). Its Avg.\
($77.19$) is within 0.11 of GIFT and 0.41 of LoRA-Loop. These results need
no generator: the 1.66\,GB buffer is about half the size of the diffusion model
in fp16 (3.4\,GB). Transfer ($69.02$) is 0.92 above ZSCL but 0.28 and 0.78
below GIFT and LoRA-Loop, so we claim higher retention at lower memory, not the
best result on every metric.

On Order~II, Last again leads ($85.74$ vs.\ $85.3$ and $85.5$), while Avg.\
($75.27$) trails by 0.4--0.6. Transfer ($63.57$) is below ZSCL ($64.2$), GIFT
($65.9$), and LoRA-Loop ($66.3$). The retention advantage therefore holds across
both orders, but the Transfer gap grows on Order~II.

\paragraph{Components.}
\Cref{tab:replay} separates the components. Without replay, distillation improves
on WE + L2 on all three metrics (+1.61 Transfer, +1.67 Avg., +1.06 Last). Adding
replay raises Avg.\ by 1.49 and Last by 2.42, to 86.87. Without $\Lzscl$, replay
alone reaches 86.62 Last but only 66.52 Transfer, and adding $\Lrd$ raises
Transfer by 1.96 for a 0.38 drop in Last. We interpret these effects
in~\cref{sec:disc-q1}.

\subsection{Replay budget (\Q{2})}
\label{sec:res-q2a}
Our hypothesis that replay saturates at a small number of exemplars per class is
false. Last rises with budget up to the largest buffer we test
(\cref{fig:q2a}), from 81.18 at 1k images to 86.17 at 11k, a gain of
4.99 from about 0.8 to 9.2 exemplars per class. \TODO{Give the slope of the
last doubling from Fig.~1, and the Transfer trend.} Details are
in~\cref{sec:abl-budget}.

\subsection{Selection (\Q{2}, \Q{3})}
\label{sec:res-select}
\TODO{After P3 and P4. One sentence for images and one for embeddings: does any
strategy beat random, and at which budget? Details in~\cref{sec:abl-select}.}

\subsection{Quantization (\Q{3})}
\label{sec:res-quant}
\TODO{Needs the storage table, which is not in this version. State how far
precision can drop before Last falls. Full sweep in~\cref{sec:abl-quant}.}

%%%%%%%%%%%%%%%%%%%%%%%%%%%%%%%%%%%%%%%%%%%%%%%%%%%%%%%%%%%%%%%%%%%%%%%%%
\section{Ablations}
\label{sec:ablations}

\subsection{Budget and distillation (\Q{2})}
\label{sec:abl-budget}
We sweep the budget $\budget$ with and without $\Lrd$, keeping all other
hyperparameters fixed. Last increases monotonically for both
variants, so buffer size is the main driver of retention regardless of $\Lrd$.
\TODO{Rewrite the rest from the sweep under the current metric: the Transfer
difference and Last cost of $\Lrd$ at each budget, and whether distillation matters
more when the buffer is small. The previous values ($+0.37$, $-1.55\rightarrow-0.27$,
the $B{=}3{,}000$ crossover) do not match \cref{tab:replay}.}

\subsection{Selection strategies (\Q{2}, \Q{3})}
\label{sec:abl-select}
\TODO{After P3 and P4: compare strategies at 1k and 11k, for images and
embeddings (\cref{tab:selection}).}

\paragraph{Allocation.}
Splitting the budget equally per task or in proportion to class count gives the
same Last (\CHECK{86.44}) and changes Avg.\ and Transfer by less than 0.10. We use
proportional allocation for its even per-class coverage, but the headline results
do not depend on it.

\subsection{Quantization (\Q{3})}
\label{sec:abl-quant}
\TODO{Embedding precision sweep (fp32, fp16, int8, int4): memory, Transfer, Avg.,
Last. Drop the PQ and reconstruction-error text if tokens stay cut.}

%%%%%%%%%%%%%%%%%%%%%%%%%%%%%%%%%%%%%%%%%%%%%%%%%%%%%%%%%%%%%%%%%%%%%%%%%
\section{Discussion}
\label{sec:discussion}

\subsection{What replay adds to distillation (\Q{1})}
\label{sec:disc-q1}

\paragraph{Replay drives retention.}
Distillation alone reaches 84.45 Last, and real exemplars add 2.42 on top of it.
Reference-data distillation therefore does not make replay redundant. This
matches how the two terms act: $\Lzscl$ and $\Lrd$ use fixed teacher text
embeddings and constrain only the image encoder, while $\Lrep$ is the only term
that updates the text encoder for previous classes. The binding constraint is
forgetting rather than plasticity, because replay raises Last by protecting
earlier tasks. \TODO{Support with BWT or the diagonal under the current metric.
The old BWT values should not be reused.}

\paragraph{Distillation preserves transfer.}
Replay alone has the lowest Transfer among the replay rows (66.52). Both anchors
raise it, but they are not equally necessary for retention: replay alone already
gives 86.62 Last, so their main role is protecting zero-shot ability.

\paragraph{$\Lrd$ can stand in for reference images.}
With no reference images, $\Lrd$ adds 1.96 Transfer to replay at a cost of 0.38
Last. It still uses Conceptual Captions text anchors, so it removes the need for
reference images but not for reference text. \CHECK{On top of $\Lzscl$, $\Lrd$
raises Transfer by about 1.7 and leaves Last essentially unchanged.}
% From the whole-number seed-42 values. Replace with exact values from the two runs.

\subsection{Quantity vs.\ fidelity (\Q{3})}
\label{sec:disc-fidelity}
Storing the final embedding (\CHECK{fp16}, 1\,kB) instead of the image (150\,kB)
reduces memory by $147\times$, from 1.66\,GB to 11\,MB at 11k samples, and costs
1.42 Last (84.75 vs.\ 86.17). Stored embeddings go stale as the image
encoder changes and send gradient only to the text encoder, so some loss is
expected. Under tight memory, however, quantity wins: 11k embeddings (11\,MB) beat
1k images (151\,MB) by 3.57 Last while using $13\times$ less memory. Fidelity
matters most when memory is plentiful, and sample count matters most when it is
scarce. \TODO{Add the equal-memory point and crossover from Fig.~2.}

\subsection{Answers and recommendations}
\label{sec:answers}
\textbf{\Q{1}:} Replay and distillation are complementary. Replay provides
retention and distillation preserves transfer, so both should be used.
\textbf{\Q{2}:} Store as much as the budget allows, since Last had not saturated
at 11k images. \TODO{Selection recommendation after P3 and P4.}
\textbf{\Q{3}:} Store images when memory allows and embeddings when it does not.
\TODO{Recommended precision and crossover point.}

\paragraph{Limitations.}
We study one benchmark and one backbone. Most runs use Order~I. Transfer also remains below the generative
methods on both orders.

\paragraph{Future directions.}
The gap between embeddings and images suggests several follow-ups: feature-level
augmentation, refreshing stored features as the encoder changes, a previous-task
teacher for feature distillation, and testing other backbones.

%%%%%%%%%%%%%%%%%%%%%%%%%%%%%%%%%%%%%%%%%%%%%%%%%%%%%%%%%%%%%%%%%%%%%%%%%
\section{Conclusion}
\label{sec:conclusion}
We studied the replay buffer in continual learning of CLIP, measuring memory in
bytes. Real exemplars add retention that distillation alone does not provide, and
with distillation they exceed the Last accuracy of generative replay without
keeping a generator at training time. Retention kept improving up to the largest
budget we tested. At tight memory budgets, embeddings are a better use of bytes
than a few images, although images remain best when memory allows. Treating the
buffer as a design choice, rather than a fixed hyperparameter, gives a simple and
strong baseline for continual vision-language learning.
